\section{Methodology}
\label{sec:method}
As shown in Figure~\ref{fig:framework}, \model adopts a plan-and-execute multi-agent architecture~\cite{huang2024understanding}, where a manager dynamically decomposes time series tasks into specialized stages executed by core agents. 
A processor then performs task-aware adaptive planning and enables parallel execution among agents to improve inference efficiency. 
For each parallel batch, the workflow consists of three agents: \ding{182} \emph{Analyzer}, which extracts statistical features $\mathbf{F}_{\text{stat}}$ and a task-aware plot image representation; \ding{183} \emph{Reasoner}, which leverages a VLM to generate anchors $\mathcal{A}$ (visual reasoning stage) and reconstruct the sequence in a latent space (numeric reasoning stage); and \ding{184} \emph{Executor}, which produces and verifies final predictions $\hat{\mathbf{Y}}$ through adaptive tool usage.

\subsection{Problem Formulation}
Let $\mathbf{X}=\{x_1,\dots,x_L\}\in\mathbb{R}^{L\times D}$ denote an input multivariate time series, where $L$ is the length of the observed sequence and $D$ is the number of variables (features). We consider four general tasks:
\ding{182}\emph{Forecasting} to predict future values $\mathbf{Y}=\{x_{L+1},\dots,x_{L+H}\}\in\mathbb{R}^{H\times D}$ over a forecasting horizon of length $H$;
\ding{183}\emph{Classification} to assign a categorical label $y\in\{1,\dots,C\}$ to the entire sequence, where $C$ is the number of classes;
\ding{184}\emph{Imputation} to estimate missing entries $\{\hat{x}_{i,j}:(i,j)\in\Omega\}$, with $\Omega$ denoting the index set of missing time-variable pairs; and
\ding{185}\emph{Anomaly Detection} to identify anomalous timestamps $\mathbf{s}\subset\{1,\dots,L\}$ within the observed window.
\model provides a unified agent-driven framework that shares common perception, reasoning, and memory components across these tasks.

\subsection{Multi-Agent Collaboration Framework}
\label{subsec:multi_agent}
Time series analysis requires heterogeneous reasoning that is difficult to unify in a single model, motivating our design of specialized agents that collaborate through shared memory to enable selective knowledge transfer while preserving task-specific representations.


\paragraph{Shared Memory Mechanism.}
We define a set of agents $\{\text{Agent}^{(i)}\}_{i=1}^{N}$ for batch data processing, where in our instantiation $N{=}3$ corresponds to \emph{Analyzer}, \emph{Reasoner}, and \emph{Executor}.
Each agent has (i) an embedding function $\text{Embed}^{(i)}(\cdot)$, (ii) a local memory view $\mathcal{M}^{(i)}$, and (iii) a communication module $\text{Com}(\cdot)$.
Cross-agent coordination is achieved through a shared memory matrix $\mathbf{M} \in \mathbb{R}^{L \times d_m}$, which integrates knowledge across agents via gated fusion.
Given an input time series $\mathbf{X}$ and current shared memory $\mathbf{M}$, each agent computes:
\begin{align}
\mathbf{h}^{(i)},\alpha_i &= \text{Agent}^{(i)}\big(\text{Embed}^{(i)}(\mathbf{X}), \mathcal{M}^{(i)}\big),\\
\text{Com}(\mathbf{h}^{(i)}|\mathcal{M}^{(i)}) &= \mathbf{h}^{(i)} + \text{LN}\Big(\text{Enc}_{\ell}(\mathbf{h}^{(i)}, \mathcal{M}^{(i)})\Big),\\
\mathbf{M} &\leftarrow \mathbf{M} \oplus \alpha_i \cdot \text{Com}(\mathbf{h}^{(i)}|\mathcal{M}^{(i)}),
\end{align}
where $\mathbf{h}^{(i)}$ denotes the agent-specific hidden state, $\alpha_i\in[0,1]$ is the confidence score controlling the update magnitude, and $\text{LN}(\cdot)$ denotes the layer normalization for $\ell =1,...,L_{enc}$ encoder layers conditioned on the agent view.
The gating mechanism $\oplus$ enables iterative memory refinement while reducing the impact of noisy agent outputs.

\paragraph{Information Carrier Encoder.}
Agents communicate through structured information carriers $\{\mathcal{I}, \mathcal{E}, \mathcal{P}\}$ that are first generated by the \textbf{Analyzer}, detailed as follows:
\begin{itemize}[leftmargin=*, itemsep=0pt]
\item $\mathcal{I}$: Plot images rendered directly from the raw time series via a rule-based operator $\mathcal{I} = \mathcal{R}(\{(t, x_{t,d})\}_{t=1}^{L})$;
\item $\mathcal{E}$: Normalized embeddings from tool-integrated time series preprocessing and encoding;
\item $\mathcal{P}$: Structured priors including statistical features, semantic descriptions, and shared memory from communication.
\end{itemize}
These carriers form a compact yet expressive interface for inter-agent communication and downstream reasoning.

\subsection{Latent Visual Reasoning}
\label{subsec:visual_reasoning}
The \textbf{Reasoner} aims to bridge intuitive visual pattern recognition and precise numerical inference for time series analysis.
Rather than directly fitting historical trajectories, we leverage visual reasoning to identify sparse but informative structural cues and subsequently enforce them through continuous-time latent dynamics.

\paragraph{VLM-based Anchoring.}
Many salient temporal structures in time series, including trends, regime shifts, and boundary behaviors, are often easier to identify from visual representations than from raw 1-D sequences.
Vision-Language Models (VLMs) perform high-level pattern recognition directly on plot images.
Formally, given plot images $\mathcal{I}$ and structured priors $\mathcal{P}$, we query the VLM as:
\begin{align}
\mathcal{A} &= \text{VLM}\Big(
\text{Enc}_{\text{vis}}(\mathcal{I}),\;
\text{Enc}_{\text{text}}(\text{Prompt}(\mathcal{P}))
\Big) \\ &= \{(t_k, v_k, \tau_k)\}_{k=1}^{K},
\end{align}
where the prompt (see Appendix~\ref{appx:prompt}) specifies the required output schema.
Each tuple consists of a critical time index $t_k \in \{1,\dots,L\}$, an estimated value $v_k \in \mathbb{R}$, and an indicator $\tau_k$ (default: $\{-1,0,1\}$) denoting local temporal behavior (e.g., falling/stable/rising).
These sparse anchors encode high-level structural knowledge extracted through visual reasoning, serving as semantic constraints that guide the latent reconstruction.

\paragraph{Latent Construction.}
While visual anchors are sparse and discrete, downstream time series tasks require dense and temporally coherent representations.
To bridge this gap, we introduce a \emph{latent construction} module that reconstructs a continuous latent trajectory by treating anchors as conditioning signals that guide latent-space evolution.
Formally, we define a continuous-time latent state $\mathbf{z}(t)$ and parameterize its dynamics by:
\begin{equation}
\frac{d\mathbf{z}(t)}{dt}
= f_{\phi}\big(\mathbf{z}(t), t \mid \mathcal{A}, \mathcal{P}\big),
\quad \mathbf{z}(t_0) = \mathbf{z}_0,
\end{equation}
where $f_{\phi}$ is a learnable dynamics function conditioned on the anchor set $\mathcal{A}$ and priors $\mathcal{P}$, and the initial state $\mathbf{z}_0$ is obtained from the embedding $\mathcal{E}$. We employ a finite-step latent reconstruction operator $\mathcal{S}(\cdot)$ to obtain latent states
$\{\mathbf{z}_t\}_{t=1}^{L+H}$, covering the observed window and the horizon:
\begin{equation}
\mathbf{u}^{(k)} = \mathcal{S}^{(k)}\!\big(\cdots\mathcal{S}^{(1)}(\mathbf{u}^{(0)}, \mathcal{A}, \mathcal{P})\big), \quad \mathbf{z}_t = [\mathbf{u}^{(k)}]_t,
\end{equation}
where $\mathbf{u}^{(i)} \in \mathbb{R}^{(L+H)\times D}$ denotes the \emph{sequence-level} latent representation after the $i$-th solver step, with $\mathbf{u}^{(0)}=\text{Proj}(\mathcal{E})$ initialized from the data embedding and fused by communication, and $\mathbf{z}_t$ being the $t$-th row of $\mathbf{u}^{(k)}$.

In practice, we instantiate $\mathcal{S}(\cdot)$ using a finite-step ($k \leq 20$) explicit latent ODE solver, as anchor conditioning substantially restricts the feasible solution space and reduces the effective curvature of the latent trajectory.
Intuitively, and consistent with prior work on continuous latent modeling~\cite{dormand1980family,rubanova2019latent},
anchor-conditioned latent dynamics act as a form of \emph{soft boundary regularization}, biasing the reconstructed trajectory toward semantically meaningful temporal structures while preserving flexibility in local variations.

\paragraph{Latent Fusion with Data Embedding.}
The reconstructed latent trajectory provides high-level structural guidance but does not replace the original data representation.
Instead, it is fused with the data embedding to form a unified latent representation using attention fusion:
\begin{equation}
\mathbf{Z} = \text{Attn}(\mathcal{E} \;\oplus\; \text{Proj}\!\left( \mathbf{u}^{(k)} \right)),
\end{equation}
where $\oplus$ denotes a gated fusion operator.
The resulting latent representation $\mathbf{Z}$ is dense, differentiable, and semantically grounded, serving as a bridge between training-free visual reasoning and data-driven numerical inference within a unified latent framework.

\subsection{Tool-Integrated Execution and Optimization}
\label{sec:execution_optimization}
The \textbf{Executor} converts latent representations into task-specific outputs via tool-integrated execution and result verification, while optimizing only coordination-critical components and keeping the reasoning modules training-free.
Instead of relying on a single monolithic predictor, \model dynamically composes a chain-of-tools that applies specialized time-series operators to refine intermediate results.

\paragraph{Tool Library for Time Series.}
At execution time, an adaptive router selects a tool chain from a predefined library $\mathcal{T}=\{\mathcal{T}_1,\dots,\mathcal{T}_M\}$, where each tool implements a specialized operation (e.g., trend extrapolation, decomposition, smoothing, statistical prediction, constraint projection, and residual correction).
Tool routing adapts to the task type $\kappa$, structured priors $\mathcal{P}$, and the shared memory $\mathbf{M}$ maintained in earlier stages. 
Formally, the executor produces a routing policy $\pi_{\Theta}$ over tools and constructs a tool sequence $\mathbf{m}=(m_1,\dots,m_J)$ accordingly. The final output is obtained by sequential composition:
\begin{equation}
\pi_{\Theta}(m \mid \kappa, \mathcal{P}, \mathbf{M})
= \mathrm{Softmax}\!\left(g_{\psi}(\kappa, \mathcal{P}, \mathbf{M})\right),
\end{equation}
\begin{equation}
\hat{\mathbf{Y}}
= \mathcal{T}_{m_J} \circ \cdots \circ \mathcal{T}_{m_1}(\mathbf{Z}),
\end{equation}
where $g_{\psi}$ is a gated controller network, and the tool sequence with length $J$ is dynamically determined by policy $\pi_{\Theta}$ and connected by $\circ$ operation.
This design enhances raw predictions through plug-and-play expert operators without retraining a monolithic backbone for a specific task.

\paragraph{Result Verifier.}
To ensure numerical validity and task consistency, \model employs a lightweight verifier that performs:
(i) format and shape correction,
(ii) constraint validation, and
(iii) task-specific normalization.
If the candidate output fails verification, the executor either recomputes or falls back to a conservative basic tool for a valid prediction.
Given a candidate output $\hat{\mathbf{Y}}$, the verifier produces the final prediction:
\begin{equation}
\mathbf{Y}^{*} = \text{Verifier}(\hat{\mathbf{Y}} \mid \mathcal{A}, \mathcal{P}),
\end{equation}
where visual anchors $\mathcal{A}$ and structured priors $\mathcal{P}$ serve as external consistency checks.
For continuous outputs, violations of anchor-implied bounds are softly corrected via constraint-aware projection; for discrete or categorical tasks, the verifier enforces schema consistency and removes logically invalid predictions.
This decouples high-level reasoning from strict constraint satisfaction, improving robustness without sacrificing flexibility.

\paragraph{Optimization Strategy.}
\model supports end-to-end differentiable execution while selectively training only coordination-critical components.
Static preprocessing, VLM prompting/inference, and selected tools in $\mathcal{T}$ are kept \emph{training-free}, serving as fixed reasoning and execution primitives.
Optimization focuses on learnable components that govern information flow and coordination, including:
(i) embedding/communication functions,
(ii) shared memory representations,
and (iii) gating and routing parameters $\Theta$.
We optimize a task-dependent objective:
\begin{equation}
\mathcal{L}
= \mathcal{L}_{\text{task}}(\mathbf{Y}^{*}, \mathbf{Y})
+ \lambda \cdot \mathcal{L}_{\text{com}}(\mathbf{M}),
\end{equation}
where $\mathcal{L}_{\text{task}}$ corresponds to the downstream objective (e.g., $\mathcal{L}_{\mathrm{MSE}}$ for forecasting),
and $\mathcal{L}_{\text{com}}$ regularizes memory updates to encourage stable inter-agent communication.

\model integrates multi-agent collaboration, intuitive visual reasoning, and tool-augmented execution into a unified framework.
By explicitly separating structural reasoning, continuous latent reconstruction, and tool-integrated execution, \model achieves strong adaptability across diverse time series tasks while preserving numerical consistency and training efficiency.
